% !TeX root = ../../protocolo.tex

This chapter presents advances in the study of the optical response of plasmonic metasurfaces composed of disordered ensembles of spherical gold nanoparticles (AuNPs), as well as the physical mechanisms that enable these systems to be used in biosensing applications. The results include both theoretical calculations aimed at the design of metasurfaces based on a specific spectral response and experimental data supporting the presence of a resonance identified in the theoretical spectra. In Section \ref{s:TDarkness}, two metasurface designs are proposed, along with their expected spectral behavior as predicted by the Coherent Scattering Model and Thin Film Theory. Additionally, the use of ellipsometric parameters ---instead of conventional reflectivity spectra--- is discussed, motivated by the phenomenon of topological darkness and the ease of performing such measurements.  Lastly, the expected spectral shift in the optical response of the metasurface due to variations on a probe medium were determined, which predict a blueshift of the identified resonance of the metasurface. To verify this phenomenon experimentally, in Section \ref{s:ExpBS} the reflectivity spectra was measured for a metasurface exposed to a probe medium artificially varying  its refractive index.